\subsection{正规收敛积分}

设 \((\mathbf{f}_{\alpha})_{\alpha\in\mathrm{A}}\) 是任意区间 \(I\subset R\) 上的一个规则函数族，取值于 R 上的完备赋范空间 E 中。假设存在 I 上的一个有限实规则函数 g，使得对每个 \(x\in I\) 与每个 \(\alpha\in A\)，\(\|f_{\alpha}(x)\|\leqslant g(x)\)，并且积分 \(\int_{I}g(t)dt\) 收敛。在这些条件之下，积分 \(\int_{I}\mathbf{f}_{\alpha}(t)dt\) 在 A 上绝对且一致收敛；事实上，对含于 I 的每个紧区间 K，

\[
\left\| \int_ {K} \mathbf {f} _ {\alpha} (t) d t \right\| \leqslant \int_ {K} g (t) d t
\]

而积分 \(\int_{I} g(t) dt\) 的收敛性蕴含：对每个 \(\varepsilon > 0\)，存在一个紧区间 \(J \subset I\)，使得对每个与 J 不相交的紧区间 \(K \subset I\)，有 \(\int_{K} g(t) dt \leqslant \varepsilon\)。当存在一个具有上述性质的实函数 g 时，我们就说积分 \(\int_{I} f_{\alpha}(t) dt\) 在 A 上正规收敛（参见 Gen.~Top., X, p.~296）。

一个积分可以在 A 上一致收敛而非正规收敛。*实函数序列 \((f_n)\) 由下列条件定义：\(f_n(x) = 1 / x\)（当 \(n \leqslant x \leqslant n + 1\)），\(f_n(x) = 0\)（当 \(x\) 在 \(\mathbf{I} = [\mathbf{0}, +\infty[\) 中取其他值），对它就出现这种情形。立即可以看出，积分 \(\int_{1}^{\infty} f_n(t) dt\) 一致收敛，但非正规收敛，因为关系式 \(g(x) \geqslant f_n(x)\)（对每个 \(x \in \mathbf{I}\) 与所有 \(n\)）蕴含 \(g(x) \geqslant 1 / x\)，从而 \(g\) 在 \(\mathbf{I}\) 上的积分不收敛。*

特别地，考虑一个级数，其通项 \(u_{n}\) 是区间 I 上的规则函数，并假设以 \(\|u_{n}(x)\|\)（它是 I 上的规则函数）为通项的级数在含于 I 的每个紧区间上一致收敛，且以 \(\int_{I}\|\mathbf{u}_{n}(t)\|dt\) 为通项的级数收敛；于是（II, p.~66, 命题 2）（规则的）函数 \(g(x)\)——即以 \(\|u_{n}(x)\|\) 为通项的级数之和——使得积分 \(\int_{I}g(t)dt\) 收敛。若令 \(f_{n}=\sum_{p=1}^{n}u_{p}\)，则积分 \(\int_{I}f_{n}(t)dt\) 正规收敛，因为有

\[
\left\| \mathbf {f} _ {n} (x) \right\| \leqslant \sum_ {p = 1} ^ {n} \left\| \mathbf {u} _ {p} (x) \right\| \leqslant g (x)
\]

对一切 \(x \in \mathrm{I}\) 与一切 \(n\) 成立；因此，和 \(\mathbf{f}\)——即以 \(\mathbf{u}_n\) 为通项的级数之和——是 \(\mathrm{I}\) 上的规则函数，且积分 \(\int_{\mathrm{I}} \mathbf{f}(t) dt\) 收敛，并有

\[
\int_ {\mathrm{I}} \mathbf {f} (t) d t = \sum_ {n = 1} ^ {\infty} \int_ {\mathrm{I}} \mathbf {u} _ {n} (t) d t\tag{11}
\]

（“非紧区间上级数的逐项积分”）。

